\section{Future work}
\label{sec:future_work}

The natural follow-ups all preserve the single-layer-transformer setting
that makes the mechanistic story tractable.

\paragraph{More group structures.}
The pattern we report --- one Fourier basis per cyclic group ---
predicts that adding a fourth task on a third cyclic group (e.g.\ a
permutation group $S_n$) should produce a third dedicated basis, and
that the multi-task model's residual stream capacity should bound how
many groups can be jointly grokked.

\paragraph{Continuous interpolation between groups.}
A small modification of the dataset --- replacing $(a, b) \mapsto (a + b)$
with $(a, b) \mapsto (a + \alpha b)$ for some smoothly varying
coefficient $\alpha$ --- would let us trace how Fourier features
transform under continuous deformations of the underlying group action.

\paragraph{Subspace overlap rather than spectrum overlap.}
Replace the cosine-similarity feature-overlap metric of Figure
\ref{fig:fig8_feature_overlap} with the principal-angle distance
between the subspaces spanned by the multi-task and single-task models'
dominant Fourier features. This would distinguish ``same frequencies in
the same proportions'' from ``same frequencies pointing in the same
embedding directions.''

\paragraph{Pruning and compression.}
The clean separation between additive and multiplicative subspaces
suggests that the multi-task model can be losslessly decomposed into
the two single-task models via an appropriate weight reparametrisation.
A successful decomposition would double as evidence that the
representations are genuinely disjoint and not merely uncorrelated on
the test set.

\paragraph{Larger primes and longer training.}
Push to $p \in \{257, 503\}$ to check whether the per-task grok-step
scaling is a fixed multiple of the single-task scaling, or whether
multi-task interference grows with $p$.
